%==============================================================================
% Research Methods - Group assignment - proposals
%==============================================================================

% This loads the article template. For a text written in English, add the
% "english" option, i.e. \documentclass[english]{hogent-article}
\documentclass{hogent-article}

% Preamble
% package imports:
\usepackage{comment} % comment-blocks
\usepackage{enumitem} % meer enum-list opties
\usepackage{tcolorbox} % deelvragen
\usepackage{todonotes} % custom todos
\usepackage{graphicx}  % nodig voor \resizebox in \fitwidth
\usepackage{ragged2e} % tekst rechts-uitlijnen (2.2)

\newcommand{\fitwidth}[1]{%
    \resizebox{\ifdim\width>\linewidth\linewidth\else\width\fi}{!}{#1}}

% Include bibliography file
\addbibresource{references.bib}

% Info about the assignment, course and study programme

\studyprogramme{Professionele bachelor toegepaste informatica}
% TODO: replace 2N1 with the correct class group(s)!
\course{Research Methods, 2AO} 
\assignmenttype{Groepsopdracht}
\academicyear{2025-2026}
\title{Theoretische ontwikkeling van een modulaire, contextbewuste Linux-desktopomgeving voor verschillende gebruikersworkflows.}

% TODO: English text? Comment or remove the lines above and uncomment the lines
% below

% \studyprogramme{Professional bachelor in applied informatics}
% \course{Research Methods, INT-IT/2B}
% \assignmenttype{Group assignment}
% \academicyear{2025-2026}
% \title{Written preparation of research proposals}

% TODO: Replace student names and email addresses
\author{Manu Coppens}
\email{manu.coppens@student.hogent.be}

\author{Peter Cypers}
\email{peter.cypers@student.hogent.be}

\author{Ewout Van Mullem}
\email{ewout.vanmullem@student.hogent.be}

\author{Stef Simons}
\email{stef.simons@student.hogent.be}

% TODO: Link your Github-repository here
\projectrepo{https://github.com/HoGentTIN/rm-25-26-groep-85}

% TODO: Choose the appropriate specialization from the list below, 
%
% - Mobile \& Enterprise development
% - AI \& Data Engineering
% - Functional \& Business Analysis
% - System \& Network Engineering
% - Mainframe Expert
%
% - If your research does not fit into any of these domains, specify a specialization yourself.
%
\specialisation{System Software \& Software Architecture}

% TODO: choose some appropriate keywords for your research proposals
\keywords{Linux, Operating Systems, Modulaire architectuur, Context-aware computing, Desktopomgeving, Resource management, Gebruikersmodi}

\begin{document}

% TODO: Replace the dummy text below with the actual contents of your research
% proposal. **ALSO REPLACE THE HEADING TITLES BY APPROPRIATE ONES!** The
% current headings are just placeholders.
\begin{abstract}
  Moderne computers worden gebruikt voor zeer uiteenlopende activiteiten zoals gamen, sofwareontwikkeling en contentcreatie. Hiervoor bestaan gespecialiseerde Linux-distributies die elk een bepaald gebruiksscenario of doelgroep centraal stellen. Pop!\_OS richt zich bijvoorbeeld sterk op desktopgebruik en ontwikkeling, terwijl SteamOS geoptimaliseerd is voor gaming. Hierdoor kan een gebruiker geconfronteerd worden met verschillende systeemomgevingen of handmatige configuraties wanneer hij regelmatig tussen meerdere workflows wisselt.
  
  Dit onderzoek analyseert de haalbaarheid van een modulaire Linux-desktopomgeving waarin verschillende gebruikersmodi, zoals development en gaming, als afzonderlijke modules kunnen worden beheerd en geactiveerd. De literatuurstudie richt zich op de architectuur van Linux-besturingssystemen, modulaire softwarearchitecturen, context-aware computing en Human-Computer Interaction. Op basis hiervan wordt een architectuur ontworpen en als proof-of-concept geïmplementeerd bovenop een bestaande Linux-distributie.
  
  Het prototype moet kunnen schakelen tussen verschillende gebruikersmodi en daarbij onder andere de desktopomgeving en systeemresources aanpassen. De evaluatie onderzoekt vervolgens in welke mate deze aanpak technisch haalbaar is en welke invloed de verschillende modi hebben op prestaties en gebruiksgemak. Het uiteindelijke doel is het onderzoeken van een modulaire desktoparchitectuur die verschillende gespecialiseerde gebruiksomgevingen binnen één systeem kan combineren.
\end{abstract}

\tableofcontents

\section{Inleiding}%
\label{sec:inleiding}

\begin{comment}
    verwachtte structuur:
    - Lead-in (algemene kadering)
    - Probleemstelling
    - Onderzoeksvraag (CoV)
    - Deelvragen probleem
    - Deelvragen oplossingen
    - doel en begrenzing vh onderzoek
    - Methode (+outline)
\end{comment}

% Lead-in:
Een moderne desktopcomputer wordt gebruikt voor zeer verschillende taken. Een gebruiker kan dezelfde computer bijvoorbeeld gebruiken voor softwareontwikkeling, gaming, studie en creatieve toepassingen. Deze activiteiten stellen echter verschillende eisen aan de softwareomgeving en aan de beschikbare systeemresources.

Binnen het Linux-ecosysteem bestaan daarom verschillende distributies die verschillende ontwerpkeuzes en gebruiksscenario's centraal stellen. Pop!\_OS, SteamOS en andere distributies gebruiken dezelfde Linux-kernel als basis, maar verschillen onder meer in software, configuratie en gebruikerservaring. Hierdoor ontstaan gespecialiseerde omgevingen voor verschillende doeleinden.

% Probleemstelling:
Hoewel Linux-distributies mogelijkheden bieden om een systeem voor een bepaald gebruiksscenario te optimaliseren, worden deze optimalisaties doorgaans gekoppeld aan een specifieke distributie of configuratie. Een gebruiker die één computer voor verschillende workflows gebruikt, moet daardoor vaak zelf applicaties, instellingen en systeemconfiguraties aanpassen aan de activiteit die hij uitvoert.

Onderzoek naar modulaire operating-system\-architecturen toont aan dat systeemfunctionaliteit kan worden opgedeeld in afzonderlijke modules en interfaces. Hierdoor kan de verdeling van functionaliteit tussen modules als een configuratiekeuze worden behandeld \autocite{}. Daarnaast is onderzocht hoe duidelijke modulegrenzen kunnen bijdragen aan de betrouwbaarheid en onderhoudbaarheid van operating systems \autocite{}. Onderzoek naar context-aware computing onderzoekt dan weer hoe systemen contextuele informatie kunnen gebruiken om hun gedrag aan te passen.

Dit onderzoek analyseerd daarom of deze principes gecombineerd kunnen worden in een modulaire Linux-desktopomgeving waarin verschillende gebruikscontexten als afzonderlijke gebruikersmodi worden aangeboden.

% Centraal onderzoeksvraag
De kernvraag van dit onderzoek luidt als volgt:
\begin{tcolorbox}[colback=yellow!10]
    Hoe kan een modulaire Linux-desktoparchitectuur ontworpen en geïmplementeerd worden die verschillende gebruikersworkflows ondersteunt door middel van contextbewuste gebruikersmodi?
\end{tcolorbox}

% Deelvragen
Hieronder wordt de onderzoeksvraag opgesplitst in een aantal deelvragen die zullen worden onderzocht om het probleem beter te begrijpen:

\begin{enumerate}
    \item Hoe zijn Linux-distributies technisch opgebouwd en welke onderdelen beïnvloeden de desktopomgeving en gebruikerservaring?
    \item Welke ontwerpkeuzes maken bestaande Linux-distributies voor verschillende gebruiksscenario's, zoals gaming, development en algemeen desktopgebruik?
    \item Welke voordelen en uitdagingen bieden modulaire software- en operating-systemarchitecturen?
    \item Welke technieken worden gebruikt om systemen contextbewust te maken en hun gedrag aan te passen aan veranderende contexten?
    \item Welke ontwerpprincipes uit bestaande Linux-distributies, modulaire systemen en contextbewuste systemen zijn toepasbaar op een modulaire desktopomgeving?
\end{enumerate}

De volgende deelvragen bouwen hierop verder en richten zich op de ontwikkeling en uitwerking van een geschikte oplossing:

\begin{enumerate}[start=6]
    \item Welke softwarearchitectuur is geschikt voor een modulaire Linux-desktopomgeving met verschillende gebruikersmodi?
    \item Welke contextinformatie kan gebruikt worden om gebruikersmodi te activeren of wijzigen?
    \item Hoe kunnen gebruikersmodi applicaties, processen en systeemresources verschillend beheren afhankelijk van de actieve module?
    \item Hoe kan een prototype worden ontworpen waarin gebruikersmodi als afzonderlijke, uitbreidbare modules kunnen worden toegevoegd en verwijderd?
    \item Welke invloed heeft het gebruik van verschillende gebruikersmodi op prestaties en gebruiksgemak?
\end{enumerate}

\section{Literatuurstudie}%
\label{sec:literatuurstudie}

Deze literatuurstudie onderzoekt de belangrijkste concepten rond modulaire en contextbewuste desktopomgevingen. Hierbij wordt gekeken naar besturingssystemen, modulariteit en contextbewustzijn. Deze begrippen vormen de theoretische basis voor het onderzoek naar een Linux-desktopomgeving die zich kan aanpassen aan gebruikers die regelmatig wisselen tussen verschillende computeractiviteiten.

\subsection{Besturingssystemen en Linux}
\label{ssec:besturingssytemen}

Een besturingssysteem (OS) is de software die de werking van een computer beheert en de communicatie tussen hardware en software mogelijk maakt. Het besturingssysteem beheert onder andere processen, bestanden, gebruikers en systeembronnen. Linux is een voorbeeld van een besturingssysteem en maakt gebruik van een kernel die de belangrijkste taken van het systeem uitvoert. Daarnaast bevat Linux verschillende diensten en hulpmiddelen voor onder andere netwerkbeheer, bestandsbeheer en beveiliging.

Volgens \textcite{Fox2021} maakt Linux bovendien een duidelijk onderscheid tussen user mode en privileged mode. Hierdoor worden gewone gebruikers en systeemonderdelen op verschillende niveaus uitgevoerd, wat bijdraagt aan de controle en beveiliging van het systeem.

\subsection{Modulaire software- en OS-architecturen}
\label{ssec:modulairesytemen}

Een modulair besturingssysteem is opgebouwd uit verschillende, relatief onafhankelijke modules die elk een specifieke taak uitvoeren. Volgens \textcite{Wichmann1969} communiceren deze modules met elkaar via een eenvoudige interface. Hierdoor kunnen afzonderlijke onderdelen van het besturingssysteem onafhankelijk worden ontwikkeld en getest. Een voordeel hiervan is dat wijzigingen of uitbreidingen aan een bepaalde module kunnen worden uitgevoerd zonder het volledige besturingssysteem te moeten aanpassen. Daarnaast kunnen modules verschillende toegangsrechten hebben, waardoor systeemmodules meer mogelijkheden hebben dan gebruikersmodules.

\subsection{Context-aware computing}
\label{ssec:contextawarecomputing}

Binnen een computer wordt telkens gewerkt in een bepaalde context. Afhankelijk van de taak, zal je een andere applicatie openen. Verschillende applicaties behoren tot dezelfde context. Volgens \textcite{Sezer2018} is de term \texttt{Context-aware computing} een manier van werken waarbij automatisch de omgeving en context wordt opgemerkt om deze automatisch aan te passen voor de taak. 

Een doel van \texttt{Context-aware computing} is zo veel mogelijk informatie te verzamelen van een apparaat of over een taak, zodat de omgeving automatisch kan worden aangepast zodat de gebruiker de taak tot zijn beste kunnen kan uitvoeren. 

\subsection{Human-Computer Interaction}
\label{ssec:hci}

De manier waarop een gebruiker gebruik maakt van een computer is heel erg belangrijk. \texttt{Human-computer interaction} is de studie naar de manieren hoe mensen computers gebruiken en de interacties tussen beide \autocite{Sharples1996}.

Volgens \textcite{Lee2021} wordt er door die connectie tussen de mens en de computer, veel focus geplaatst op het design van systemen en applicaties. Hoe sterker de mensen met de computers kunnen werken, hoe productiever er gewerkt kan worden.

\subsection{State of the Art}
\label{ssec:stateoftheart}

\section{Methodologie}
\label{sec:methodologie}

Dit onderzoek wordt uitgevoerd volgens de Proof-of-Concept-methodologie. Deze onderzoeksmethode combineert een theoretische analyse van een probleem met het ontwerpen, ontwikkelen en evalueren van een softwareoplossing. Het onderzoek wordt opgedeeld in vijf opeenvolgende fasen, waarbij de resultaten van elke fase de basis vormen voor de volgende. Dit wilt echter niet zeggen dat een fase volledig afgesloten wordt na het uitvoeren. Indien nodig, kunnen de vorige fases opnieuw doorlopen worden en wordt iteratief werken mogelijk.

\subsection{Fasering}%
\label{ssec:fasering}

Het onderzoek loopt in zes fasen. Tabel~\ref{tab:fasering} toont welke deelvraag
in welke fase aan bod komt.

\begin{table*}[t]
    \centering\small
    \caption{Deelvragen, fasen en methoden}
    \label{tab:fasering}
    \fitwidth{
        \begin{tabular}{@{}llp{6.5cm}p{4cm}@{}}
            \toprule
            \textbf{Fase} & \textbf{Deelvragen} & \textbf{Methode} & \textbf{Resultaat} \\
            \midrule
            1. Literatuurstudie      & 1--5       & Literatuuronderzoek               & Referentiekader \\
            2. State of the Art      & 2, 5       & Documentanalyse, vergelijking     & Vergelijkingsmatrix \\
            3. Architectuurontwerp   & 6--9       & Ontwerp, beslissingsmatrix        & Desktop-architectuur \\
            4. Implementatie         & 7--9       & Softwareontwikkeling              & Proof-of-concept \\
            5. Evaluatie en  analyse & 10         & Performance- en usabilitytests    & Meetresultaten \\
            6. Rapportering          & 1--10      & Analyse van bevindingen           & Conclusie \\
            \bottomrule
    \end{tabular}}
\end{table*}

\subsection{Fase 1 – Literatuurstudie en probleemanalyse}
\label{ssec:fase1}

In de eerste fase wordt een literatuurstudie uitgevoerd naar de belangrijkste concepten die relevant zijn voor het onderzoek. Hierbij wordt onderzocht hoe Linux-besturingssystemen en Linux-distributies zijn opgebouwd, welke principes worden gebruikt binnen modulaire software- en operating-systemarchitecturen en hoe contextbewuste systemen contextinformatie gebruiken om hun gedrag aan te passen. Daarnaast wordt onderzocht welke aspecten van Human-Computer Interaction relevant zijn voor het aanpassen van een desktopomgeving aan verschillende gebruikersactiviteiten.

Op basis van de literatuur wordt het onderzoeksprobleem verder afgebakend en worden relevante begrippen en ontwerpprincipes geïdentificeerd. De resultaten van deze fase vormen het theoretische kader voor de verdere analyse en ontwikkeling van het proof-of-concept.

\subsection{Fase 2 – State of the Art}
\label{ssec:fase2}

In de tweede fase worden bestaande Linux-distributies en relevante systemen systematisch met elkaar vergeleken. Hierbij worden onder andere Ubuntu, Pop!\_OS, SteamOS en NixOS onderzocht. De vergelijking richt zich op aspecten die relevant zijn voor het onderzoek, zoals de systeemarchitectuur, modulariteit, configuratiemogelijkheden, desktopomgeving, resourcebeheer en de manier waarop het systeem is afgestemd op specifieke gebruiksscenario's.

Daarnaast worden bestaande oplossingen voor modulaire, contextbewuste en aanpasbare systemen onderzocht. Op basis van deze vergelijking wordt bepaald welke ontwerpprincipes en technieken relevant zijn voor het voorgestelde systeem en welke beperkingen in bestaande oplossingen aanwezig zijn. De resultaten worden samengebracht in een overzicht van de huidige stand van zaken en vormen, samen met de literatuurstudie, de basis voor de requirements van het proof-of-concept.

\subsection{Fase 3 – Architectuurontwerp}
\label{ssec:fase3}

In de derde fase wordt op basis van de literatuurstudie en de State of the Art een architectuur voor het proof-of-concept ontworpen. Hierbij worden de functionele en niet-functionele vereisten vastgelegd en wordt bepaald welke systeemcomponenten noodzakelijk zijn.

Er wordt onderzocht hoe de verschillende modules worden opgebouwd en hoe deze met elkaar en met de onderliggende Linux-omgeving communiceren. Daarnaast wordt bepaald welke contextinformatie kan worden gebruikt om gebruikersmodi te activeren en hoe een actieve modus invloed kan uitoefenen op applicaties, processen en systeeminstellingen. De belangrijkste systeemcomponenten en hun onderlinge relaties worden beschreven aan de hand van architectuur- en UML-diagrammen.

Het resultaat van deze fase is een technisch ontwerp dat als basis dient voor de implementatie.

\subsection{Fase 4 – Implementatie}
\label{ssec:fase4}

Tijdens de vierde fase wordt het ontworpen systeem geïmplementeerd als een proof-of-concept bovenop een bestaande Linux-distributie. De implementatie richt zich op de kernfunctionaliteiten van het voorgestelde systeem en niet op het ontwikkelen van een volledig nieuw besturingssysteem.

De focus ligt op het realiseren van de modulaire architectuur, het implementeren van de contextdetectie en het activeren en beheren van een beperkt aantal gebruikersmodi. Voor de proof-of-concept worden minimaal twee verschillende gebruiksmodi geïmplementeerd, bijvoorbeeld Development en Gaming. Afhankelijk van de gekozen architectuur kunnen deze modi onder andere desktopinstellingen en systeemresources configureren.

De ontwikkeling gebeurt iteratief. De verschillende componenten worden afzonderlijk getest en vervolgens geïntegreerd tot één werkend systeem. Het resultaat van deze fase is een functioneel proof-of-concept.

\subsection{Fase 5 – Evaluatie en analyse}
\label{ssec:fase5}

In de vijfde fase wordt het proof-of-concept geëvalueerd aan de hand van vooraf bepaalde gebruiksscenario's en meetcriteria. Eerst wordt gecontroleerd of de verschillende functionaliteiten correct werken. Hierbij wordt onder andere nagegaan of gebruikersmodi correct kunnen worden geactiveerd, of de bijbehorende modules correct worden geladen en of de contextdetectie naar behoren functioneert.

Vervolgens wordt onderzocht welke invloed het gebruik van de verschillende modi heeft op het systeem. Hiervoor worden relevante prestatie-indicatoren, zoals CPU- en geheugengebruik en applicatie- of gameprestaties, gemeten en vergeleken met een referentiesituatie. Daarnaast wordt, binnen de beschikbare tijd en scope van het onderzoek, de bruikbaarheid van het systeem geëvalueerd aan de hand van vooraf bepaalde gebruiksscenario's.

De verzamelde resultaten worden geanalyseerd en gebruikt om de technische haalbaarheid, voordelen en beperkingen van de voorgestelde architectuur te beoordelen. Waar mogelijk worden vastgestelde problemen nog tijdens deze fase verbeterd en opnieuw geëvalueerd.

\subsection{Fase 6 – Rapportering en afronding}
\label{ssec:fase6}

In de laatste fase worden de resultaten van het onderzoek verwerkt in een rapport. De bevindingen uit de literatuurstudie, de State of the Art, het architectuurontwerp, de implementatie en de evaluatie worden samengebracht.

Op basis van de onderzoeksresultaten wordt de centrale onderzoeksvraag beantwoord. Daarnaast worden de beperkingen van het onderzoek en het proof-of-concept besproken. Hierbij wordt onder andere aangegeven welke functionaliteiten buiten de scope van het onderzoek vallen en welke onderdelen in toekomstig onderzoek verder kunnen worden ontwikkeld.

Tot slot wordt de voorbereiding van de verdediging uitgevoerd.

\subsection{Doel en begrenzing}
\label{ssec:beperkingen}

Het doel van dit onderzoek is het ontwerpen en evalueren van een proof-of-concept van een modulaire Linux-desktopomgeving met contextbewuste gebruikersmodi. Het onderzoek richt zich op de desktoplaag en de interactie met systeemprocessen en resources. Het ontwikkelen van een eigen Linux-kernel, hardwaredrivers, package manager of volledig onafhankelijke Linux-distributie valt buiten de scope.

Het prototype wordt ontwikkeld bovenop een bestaande Linux-distributie. Voor de proof-of-concept worden maximaal twee gebruikersmodi geïmplementeerd. De architectuur wordt wel ontworpen met uitbreidbaarheid als uitgangspunt, zodat bijkomende modi in de toekomst kunnen worden toegevoegd.

Automatische adaptatie op basis van machine learning of langdurige gebruikersprofilering valt eveneens buiten de scope. Contextdetectie wordt in het prototype beperkt tot vooraf bepaalde, observeerbare gebeurtenissen en/of expliciete gebruikerskeuzes.

\section{Verwacht resultaat en conclusie}
\label{sec:results}

Het verwachte resultaat is een proof-of-concept van een modulaire Linux-desktopomgeving waarin minimaal twee gebruikersmodi kunnen worden geactiveerd. Elke modus zal een eigen configuratie bevatten voor onder andere desktopinstellingen en systeemresources.

Daarnaast wordt een architectuurmodel opgesteld waarin wordt beschreven hoe modules worden beheerd, hoe contextinformatie wordt verwerkt en hoe een actieve gebruikersmodus de werking van het systeem kan beïnvloeden.

De evaluatie moet inzicht geven in de technische haalbaarheid van deze aanpak en in de mogelijke invloed van contextspecifieke configuraties op systeemprestaties en gebruiksgemak.

\printbibliography[heading=bibintoc]

\end{document}